\subsection{A peak section on the actual thick quotient}
\label{shared-bergman-actual-thick}

The following argument concerns the original good cover and its
Hermitian line, at one coefficient embedding.  It uses a complete
quotient metric, so no estimate for a separately periodized peak or for
an arithmetic fundamental volume is needed.  Write
$f=[E:\mathbb Q]/2$, $A=fr$, and $D=f\binom r2$.
At an orthogonal frame the sign matrix is
$J=\operatorname{diag}(1_p,-1_q)$, where $p,q\geq1$ and $p+q=r$.
For an integral weight $\lambda$ put
$S=\sum_{\tau,i}|\lambda_{\tau i}|$.
The original norm of a section is denoted by $\|P\|_\lambda$.

\begin{theorem}[The actual thick peak]
\label{bergman-actual-thick-peak}
There is an absolute constant $C\geq1$ with the following property.
On any component of the original good cover, suppose
\[
 \begin{gathered}
 \lambda_{\tau i}-\lambda_{\tau,i+1}\geq CA^8
       \quad(i+1\ne p),\\
 \lambda_{\tau,r-1}-\lambda_{\tau0}\geq CA^8
       \quad\text{for every }\tau.
 \end{gathered}
\]
Let $x$ be an orthogonal frame with
\begin{equation}
 \label{bergman-actual-thick-test}
  M\,\operatorname{firstMinimumSq}(x)\geq100A^4,
\end{equation}
where $M$ is the original level and the auxiliary level factor is one.
Then the original cusp space contains a section $P$ with $P(x)\ne0$ and
\begin{equation}
 \label{bergman-actual-thick-peak-bound}
 e^{-C(S+A^8)}\int_{\mathcal F}\|P\|_\lambda^2\,d\mu_F
 \leq \|P(x)\|_\lambda^2\,\mu_F(\mathcal F)
\end{equation}
for every original fundamental set $\mathcal F$.
The same conclusion holds whenever the original chart ball of radius
$1/(100r^2)$ injects in the quotient, without assumption
\eqref{bergman-actual-thick-test}.
\end{theorem}

All arithmetic hypotheses here are those of the constructed good cover:
the original self-dual lattice, full level markings, fine moduli
uniformization and its exact automorphic dictionary.  The assertion is
not an extension to an arbitrary analytic lattice with an unspecified
multiplier.  In particular it applies to the untagged weights in the
original thick Bergman consumer.\cite{bergman-thick-source}

\begin{proof}
\emph{A tensor test for a short return.}
For Euclidean lattices $L_1,L_2$ and a nonzero integral tensor
$T\in L_1\otimes L_2$ of rank $k$,
\begin{equation}
 \label{bergman-tensor-minimum}
 \|T\|_{\mathrm{HS}}\geq
      \frac{\lambda_1(L_1)\lambda_1(L_2)}{\sqrt{k}}.
\end{equation}
Indeed, choose a basis $u_1,\ldots,u_k$ of the saturated first support
in $L_1$ and write $T=\sum_i u_i\otimes v_i$.  The $v_i$ are independent
elements of $L_2$.  If $U,V$ are their Gram matrices, the arithmetic--
geometric mean inequality for the positive eigenvalues of
$U^{1/2}VU^{1/2}$ gives
\[
 \|T\|_{\mathrm{HS}}^2
 \geq k(\det U\det V)^{1/k}.
\]
The cube form of Minkowski's theorem gives
$\operatorname{covol}(L)^{1/k}\geq\lambda_1(L)/\sqrt{k}$
for every rank-$k$ Euclidean lattice.  Apply it to the two support
lattices, whose first minima are at least those of $L_1,L_2$.
This proves \eqref{bergman-tensor-minimum} without a restriction on rank.

For the original lattice $\Lambda$, use the actual majorant
\[
 Q_x(v)=\sum_\tau2|\tau\xi_0|\,
                 \operatorname{maj}_{x,\tau}(v_\tau).
\]
Its real dimension is $N=2fr$.  If $L$ is its Euclidean realization,
the original alternating trace self-duality identifies $L^*$
isometrically with $L$.  More explicitly, the two trace-dual
realizations $\phi,\psi$ satisfy
\[
 \langle\phi u,\psi w\rangle=\operatorname{Tr}_{E/\mathbb Q}h(u,w),
 \qquad \|\psi(-\xi_0w)\|^2=Q_x(w),
 \qquad L^*=\psi(-\xi_0\Lambda).
\]
Thus $\lambda_1(L)^2=\lambda_1(L^*)^2$
is exactly $\operatorname{firstMinimumSq}(x)$; there is no division by
$f$ and no replacement by a rank-one covolume.

Suppose $\gamma$ preserves $h,\Lambda$, is congruent to one modulo $M$,
and takes two points of the stated small chart ball to the same flag.
The two chart majorants differ from $Q_x$ by factors between $1/2$ and
$2$, so $\|\gamma\|_{\mathrm{op}}\leq2$ in $Q_x$.
For completeness, write a chart frame as $x(I+N(z))$, with $N$ strictly
triangular.  Each coordinate is bounded by $\|z\|$, and
$\|N\|_{\mathrm{op}}\leq r\|z\|\leq1/(100r)$.
The leading Gram matrices are perturbations of the corresponding
blocks of $J$.  Indefinite Gram--Schmidt, or its successive Schur
complements, gives a unitary representative with norm and inverse norm
at most $\sqrt2$.  Its ambiguity is the compact diagonal torus.
This proves the asserted majorant comparison and return bound.

If $\gamma\ne1$, then
$(\gamma-1)/M\in\operatorname{End}_{\mathbb Z}\Lambda=L\otimes L^*$
is nonzero.  Equation \eqref{bergman-tensor-minimum} yields
\[
 M\operatorname{firstMinimumSq}(x)
 \leq \sqrt N\,\|\gamma-1\|_{\mathrm{HS}}
 \leq3N=6A.
\]
This contradicts \eqref{bergman-actual-thick-test}.

One must still pass from this principal subgroup to the original
uniformizing group.  The actual moduli comparison supplies this step.
An identification of two flags by an original return gives an
$\mathcal O_E$-linear, polarization-preserving isomorphism of their
marked complex tori.  On their homology lattice it is a unitary
$\gamma'$ preserving $\Lambda$.  Preservation of the two original
level markings gives congruence on the two conjugate level ideals.
Their coprimality and trace self-duality give
$(\gamma'-1)\Lambda\subset M\Lambda$.  In the general mixed level the
modulus is $\operatorname{lcm}(M,A_{\rm lev})$; here
$A_{\rm lev}=1$, so it is precisely $M$.
The two lifted flags differ by the original upper triangular frame
change.  The preceding argument makes $\gamma'=1$; faithfulness of
the fine level uniformization then makes the original return trivial.
This uses the actual orbit comparison, not an assertion that its
original group is itself principal.\cite{bergman-thick-source}

\emph{A complete metric with a uniform local chart.}
Let $\mathscr D=\prod_\tau U(p,q)/T$ be the full flag domain and
$Y=\Gamma\backslash\mathscr D$ its actual component.
For one embedding put
\[
 \lambda^0_i=\begin{cases}-i&i<p,\\ r-i&i\geq p.
                        \end{cases}
\]
The original squared line metric in the chart $I+N(z)$ is
$e^{-\varphi_\lambda}$, where
\begin{equation}
 \label{bergman-exact-metric}
 \varphi_\lambda(z)=\sum_{\tau,i}\lambda_{\tau i}
       \bigl(\log|\Delta_{\tau,i+1}(z)|-
             \log|\Delta_{\tau,i}(z)|\bigr).
\end{equation}
Here $\Delta_{\tau,j}$ is the leading $j$ by $j$ Gram determinant,
and $\Delta_{\tau,0}=1$.
Differentiation at zero gives the coefficient
\[
 \epsilon_i\epsilon_j(\lambda_{\tau i}-\lambda_{\tau j})
       \quad(i<j),\qquad
 \epsilon_i=\begin{cases}1&i<p,\\-1&i\geq p.\end{cases}
\]
For $\lambda^0$ these coefficients are between $1$ and $r$.
Write $\omega_0=i\partial\bar\partial\varphi_{\lambda^0}$.
It is a positive invariant K\"ahler form.  Its homogeneous Riemannian
metric is complete, and so is the quotient metric on $Y$.
If $m$ is the least of the gaps in the theorem, telescoping the gaps,
including the gap from the last to the first weight, gives globally
\begin{equation}
 \label{bergman-curvature-positive}
 i\Theta(L_\lambda)\geq m\omega_0.
\end{equation}
The invariant canonical weight is the sum of the flag roots; its
successive coefficients have absolute value at most $2r$.
The same root calculation therefore gives
$\operatorname{Ric}(\omega_0)\geq-Cr\omega_0$.

There is an absolute $c>0$ such that, on the coordinate ball of radius
\[
 \eta=cA^{-3},
\]
the metric $\omega_0$ is between one half and $2r$ times the Euclidean
metric, and $|\varphi_\lambda|\leq CS$.
Here and below $c$ can be decreased by an absolute factor.
These are uniform matrix estimates, not an appeal to compactness in
varying dimension.  On this ball each leading Gram matrix has inverse
norm at most $2$.  Differentiate
$d\log\det G=\operatorname{tr}(G^{-1}dG)$ three times.
For unit Hilbert--Schmidt directions every resulting trace is bounded
by $Cr$; the sum of the absolute coefficients of $\lambda^0$ is at
most $r^2$.  Thus the third derivative is bounded by $Cr^3$ on each
embedding block.  The blocks are orthogonal, so this is also an
operator bound for their product.  It gives the stated Hessian
comparison for $\eta=cA^{-3}$.  The same determinant formula bounds
\eqref{bergman-exact-metric} by $CS$.
This smaller ball lies in the injective original chart.

\emph{Solve once on the quotient.}
Choose a holomorphic line frame $e$ with $\|e(x)\|=1$ and a smooth
cutoff $\chi$, equal to one for $|z|\leq\eta/8$, supported in
$|z|<\eta/4$, with $|\bar\partial\chi|_{\omega_0}\leq C/\eta$.
Let $\chi_0$ be one on $|z|\leq\eta/2$, zero for $|z|\geq\eta$,
and define on the quotient
\[
 \psi_\varepsilon=(D+1)\chi_0(z)
       \log\frac{|z|^2+\varepsilon^2}{\eta^2+\varepsilon^2},
       \qquad0<\varepsilon<\eta/16.
\]
Extend it by zero outside the chart.  Then $\psi_\varepsilon\leq0$ and
\[
 i\partial\bar\partial\psi_\varepsilon
       \geq-C(D+1)\eta^{-2}\omega_0\geq-CA^8\omega_0.
\]
The logarithm is plurisubharmonic where $\chi_0=1$; on the cutoff
annulus it and its first two scaled derivatives have absolute bounds.
Since $D\leq A^2$, the power $A^8$ includes every dimension factor.

Regard $\bar\partial(\chi e)$ as a $(D,1)$ form with values in
$L_\lambda\otimes K_Y^{-1}$.  Its canonical identification with an
$L_\lambda$-valued $(0,1)$ form is isometric for the induced metrics.
By \eqref{bergman-curvature-positive}, the Ricci estimate and a large
enough absolute $C$ in the gap hypothesis, this bundle with weight
$e^{-\psi_\varepsilon}$ has curvature at least $\omega_0$.
The complete K\"ahler $L^2$ theorem\cite{bergman-demailly}
therefore solves
$\bar\partial u_\varepsilon=\bar\partial(\chi e)$ with
\begin{equation}
 \label{bergman-solved-norm}
 \int_Y\|u_\varepsilon\|_\lambda^2e^{-\psi_\varepsilon}\,dV_0
 \leq C\eta^{-2}e^{C(S+D)}\operatorname{vol}_0(B_\eta).
\end{equation}
The right side follows on the support of $\bar\partial\chi$, where
$e^{-\psi_\varepsilon}\leq64^{D+1}$.
No pseudoconvexity assumption on the quotient is used.

Take a weak $L^2$ limit.  On every closed annulus about $x$ the weights
converge uniformly, and weak lower semicontinuity followed by monotone
exhaustion gives \eqref{bergman-solved-norm} with
$\psi=(D+1)\chi_0\log(|z|^2/\eta^2)$.
The limit $u$ is holomorphic near $x$.  Its weighted square integrability
forces $u(x)=0$.  Consequently $P=\chi e-u$ is global and holomorphic,
$P(x)=e(x)$, and
\begin{equation}
 \label{bergman-quotient-l2}
 \int_Y\|P\|_\lambda^2\,dV_0
       \leq e^{C(S+A^8)}\operatorname{vol}_0(B_\eta).
\end{equation}

\emph{Membership in the original cusp space.}
This last point requires the actual arithmetic quotient and the deep
weight range.  Restriction to a compact flag fibre gives the fixed
Borel--Weil representation of $U(p)\times U(q)$.
The resulting automorphic function is $K$-finite, and holomorphicity
gives its holomorphic discrete-series infinitesimal character.
One convention for its lowest compact highest weight is
$\mu=-w_{0,K}\lambda$.
The associated finite-dimensional highest weight is
\[
 \nu=\mu-(\underbrace{q,\ldots,q}_{p},
                   \underbrace{-p,\ldots,-p}_{q}).
\]
The adjacent gaps of $\nu$ inside the two blocks are at least $m$;
its central gap is at least $m-r$.  Thus $\nu$ is integral and strictly
dominant at every real place.  A simultaneous dual convention has the
same regularity property.

The $L^2$ holomorphic section defines a square-integrable automorphic
form in the corresponding coherent degree-zero space.
To check the growth qualification, a local holomorphic mean-value
estimate on a fixed translated chart bounds its value by its local
$L^2$ mass.  The overlap multiplicity is bounded by the number of
integral matrices in a ball of radius polynomial in
$\|g\|+\|g^{-1}\|$.  This is a polynomial count on the fixed
arithmetic lattice.  It gives moderate growth; the highest-weight
relations give finiteness under the centre of the enveloping algebra.
These are qualitative assertions at each fixed setting and introduce
no constant into \eqref{bergman-quotient-l2}.
The actual dictionary supplies the algebraic automorphic bundle.
Its multiplier on the present tag-zero weight is trivial, as follows
directly from the constructed good-cover law.  For a nonnegative
antitone exponent array $a$, that law gives
$\chi(\operatorname{laU}(a,0),\gamma)=1$.
At tag zero, $\operatorname{laU}$ is the signed permutation taking
the first block in reverse order with a minus sign and the second
block in its original order with a plus sign, with the prescribed
conjugate embedding relabeling.  Given any integral flag weight,
pull it back to an integer array $c$.  Choose an integer $B$ larger
than every $|c_j|$ and $|c_j-c_{j+1}|$, including the last jump to
zero.  The arrays $b_j=B(r-j)$ for $j<r$ and zero afterward, and
$a=b+c$, are nonnegative and antitone.  Since
$\lambda=\operatorname{laU}(a,0)-\operatorname{laU}(b,0)$,
the original weight-addition law gives $\chi(\lambda,\gamma)=1$.
Thus no inference from finite character image to a congruence kernel,
and no new level cover, is needed.

For unitary groups, square-integrable coherent cohomology at the
infinitesimal character of an integral strictly dominant weight equals
cuspidal coherent cohomology.\cite{bergman-harris}
The cited formulation uses the unitary similitude Shimura variety.
Use its connected component with a compact finite level intersecting
the unitary group in the given level and with similitude in
$\widehat{\mathbb Z}^{\times}$.  A positive rational similitude in
that finite level is one.  Consequently its component arithmetic
group is exactly the original unitary arithmetic group.
An integral unitary weight extends to the similitude group over
$\mathbb C$; its complexification is a product of the original
general linear factors and a multiplicative group.  This identifies
the same coherent bundle on the same component, with no change of
the section or its metric.  The preceding calculation checks the
regularity hypothesis of the cited result, so $P$ is cuspidal.
Rapid decay of a cuspidal automorphic form of moderate
growth, applied to this finite compact type, makes its original
invariant point norm bounded.\cite{raw-theta-rapid-decay}
The constructed good-cover dictionary therefore puts it in the
literal space $X.\mathrm{SS}(\lambda)$.
In particular, square integrability alone has not been used to infer
cuspidality.

Finally $dV_0$ and the original flag Haar measure differ by a positive
constant, since both are invariant.  The injective ball has volume at
most that of the quotient.  These two facts turn
\eqref{bergman-quotient-l2} into
\eqref{bergman-actual-thick-peak-bound}, with every original Haar
normalization intact.
\end{proof}

\paragraph{The exact Bergman consequence.}
On the actual component extend $P$ by zero to the other components.
The original dictionary identifies its invariant point norm and its
component mean square.  A component's probability mass in the original
global Petersson norm is at most one.  The variational characterization
of the squared norm of evaluation therefore gives
\[
 \operatorname{bergOn}(y)\geq e^{-C(S+A^8)}
\]
at every point satisfying \eqref{bergman-actual-thick-test}; orthogonal
frame transfer preserves both sides.  This is the exact comparison
used by the original component Bergman consumer, rather than a change
to an independently normalized $L^2$ space.
For its untagged weight $u_1$, the model permutation preserves the
absolute weight sum, so $S=\ell_1(u_1)$.
If the original numerical data give
$\operatorname{lenT}\ell_1(u_1)\leq C_2\operatorname{unitE}$ and
$\operatorname{lenT}A^8\leq\operatorname{unitE}$, then the same
$B=C(\ell_1(u_1)+A^8)$ satisfies
$\operatorname{lenT}B\leq C(C_2+1)\operatorname{unitE}$.
These numerical inequalities must be checked in their original host
range when invoking that asymptotic wrapper.
The theorem supplies the thick analytic estimate and its actual
arithmetic injectivity.  The thin estimate, global upper section-frame
bound, point-height estimate and final section remainder remain
separate obligations.
